
\subsubsection{Estrutura em malha fechada}

A técnica de controle aplicada consiste em uma estrutura em malha fechada na qual:

\begin{enumerate}
\item a dinâmica orbital relativa é descrita pelas equações de \gls{hcw}, que fornecem a evolução da posição e velocidade do centro de massa da base no referencial \gls{lvlh};
\item a atitude da base é estabilizada em torno de uma referência desejada por meio do controlador PD de atitude \eqref{eq:controle_PD_atitude}, de forma que o alinhamento da espaçonave perseguidora se mantenha em relação ao alvo;
\item as juntas do \gls{mr} seguem trajetórias de referência obtidas pela cinemática inversa, sendo rastreadas pelo controlador \gls{pd} articular \eqref{eq:controle_PD_juntas}.
\end{enumerate}

Essa combinação permite que, além da transição das dinâmicas utilizadas neste trabalho, durante as fases de aproximação curta e final, a base física mantenha uma atitude compatível com os requisitos de aproximação, enquanto o manipulador robótico realiza a tarefa de alinhamento fino e acoplamento com a porta de acoplamento.
